In this section we consider the problem of checking whether the language
recognized by a UHAT or B-RASP program is non-empty. In particular, the
technique is essentially a simulation of a Turing machine with an
$2^{O(n)}$-sized tape (for a given $n$). As we will see later, 
there are Turing machines such that
the smallest (i.e.\ shortest) accepted string 
by the constructed UHAT
is of length at least $2^{2^{\Omega(n)}}$.

\begin{example}\label{ex:tiling}
To illustrate the idea, we describe a B-RASP program that accepts strings of the form
\[0000a_1 \# 0001a_2 \# 0010a_3 \# \dots \# 1111a_{2^4} \#\]
where $a_i \in \{a,b,c\}$ such that $(a_j,a_{j+1}) \in H$ for all $1 \le j < 2^4$.
Here, $H := \{(a,b),(b,c),(b,a),(c,b)\}$ is a set of constraints specifying which symbols can be next to each other.
%Assume for simplicity that the input has already the right format, i.e., is contained in $(\{0,1\}^4\{a,b,c\}\#)^*$.
%Moreover, it is easy to see that a B-RASP program can succinctly check that the first bit sequence is 0000 and the last is 1111.
For simplicity, we concentrate on the two main conditions: (i) checking that the bit counter is incremented and 
(ii) checking that the successive symbols are in $H$.
To check (i), we use the following attention operation:
\begin{multline*}
C_{+1}(i) :=\ \maxatt_j[j<i,Q_\#(j)]\ \bigvee_{k=1}^4 \big(\bigwedge_{r=1}^{k-1} \neg C_r(i) \wedge C_r(j)\big) \wedge
C_k(i) \wedge \neg C_k(j)\ \wedge\\ \ \big(\bigwedge_{r=k+1}^{4} C_r(i) \leftrightarrow C_r(j)\big) : 1
\end{multline*}
Assume $i$ is a $\#$-position.
Attention selects the rightmost $\#$-position $j$ left of position $i$.
Let $b_1^i \dots b_4^i$ and $b_1^j \dots b_4^j$ be the bit strings directly left of position $i$ and $j$, respectively.
We assume that we already defined $C_k(i) = b_k^i$ and $C_k(j) = b_k^j$ for all $k \in [1,4]$.
Then the above value predicate checks that the binary number $b_1^i \dots b_4^i$ is the number $b_1^j \dots b_4^j$ incremented by 1.
To check (ii), we can use the attention operation
\begin{equation*}
M_{\leftarrow}(i) :=\ \maxatt_j[j<i,Q_a(j) \vee Q_b(j) \vee Q_c(j)]\ \bigvee_{(h,h') \in H} Q_h(j) \wedge Q_{h'}(i) : 1.
\end{equation*}
If $i$ is a position of a symbol $a_i$, attention picks the rightmost position $j$ of a symbol $a_j$ to the left of $i$ and 
checks with the value predicate that $(a_j,a_i) \in H$.

This allows us to succinctly recognize a language whose smallest string has length exponential 
in the number of bits of the binary counter.
By stacking multiple such strings vertically and 
introducing vertical constraints in addition to the horizontal constraints $H$,
we can even succinctly recognize languages whose smallest string has doubly exponential length. 
\end{example}

We prove the following precise complexity bounds:
\begin{theorem}\label{thm:UHAT-emptiness}
The non-emptiness problem for UHATs and B-RASP programs is $\EXPSPACE$-complete.
\end{theorem}

We start with the lower bound and show it first for B-RASP programs.
\begin{proposition}\label{prop:hardness}
The non-emptiness problem for B-RASP programs is $\EXPSPACE$-hard.
\end{proposition}
For the proof we use the techniques illustrated in \cref{ex:tiling} and reduce from a so-called \emph{tiling problem}.
A \emph{tile} is a quadruple $t \in \N_0^4$, where we write $t = \langle\tleft(t),\tup(t),\tright(t),\tdown(t)\rangle$.
The \emph{$2^n$-tiling problem} is defined as follows:

\begin{description}
\item[Given:] An integer $n > 0$ in unary, a finite set $T$ of tiles, and a tile $t_\mathit{fin} \in T$
\item[Question:] Do there exist $m > 0$ and a function $\tau \colon \{1,\dots,2^n\} \times \{1,\dots,m\} \to T$ such that
\begin{enumerate}
\item $\tau(2^n,m) = t_\mathit{fin}$, \label{itm:t-fin}
\item $\tdown(\tau(i,1)) = \tup(\tau(i,m)) = 0$ for all $1 \le i \le 2^n$, \label{itm:bot-top}
\item $\tleft(\tau(1,j)) = \tright(\tau(2^n,j)) = 0$ for all $1 \le j \le m$, \label{itm:first-last}
\item $\tright(\tau(i,j)) = \tleft(\tau(i+1,j))$ for all $1 \le i < 2^n$ and $1 \le j \le m$, and \label{itm:right-left}
\item $\tup(\tau(i,j)) = \tdown(\tau(i,j+1))$ for all $1 \le i \le 2^n$ and $1 \le j < m$? \label{itm:up-down}
\end{enumerate}
\end{description}
The following is shown in \citep{Schwarzentruber19}:
\begin{proposition}\label{prop:tiling}
The $2^n$-tiling problem is $\EXPSPACE$-complete.
\end{proposition}

The reduction to B-RASP uses an encoding of the function $\tau$ as a sequence of strings 
which are of a similar form as in \cref{ex:tiling}, but is substantially more
involved.
The key observation is that strict future masking with rightmost tie-breaking 
enables us to check conditions between successive tiles (condition~\ref{itm:right-left}) 
but also between the current tile and the tile at the most recent past occurrence of the same counter value
(condition~\ref{itm:up-down}).
The full proof can be found in \cref{app:emptiness}.

We observe that the B-RASP program constructed in the proof of \cref{prop:hardness} can easily be converted to UHAT,
which yields the $\EXPSPACE$ lower bound also for UHAT.
\begin{proposition}\label{prop:UHAT-hardness}
The non-emptiness problem for UHAT is $\EXPSPACE$-hard.
\end{proposition}
\begin{proof}[Proof sketch]
We show that the B-RASP program constructed in the proof of \cref{prop:hardness} can be converted to a UHAT in polynomial time.
To this end, note that Boolean operations can easily be simulated using affine transformations and ReLU.
For the attention operations we use an attention layer.
The value predicates $V(i,j)$ can be simulated by copying the required components of the $j$-th vector,
that was selected by attention, to the $i$-th vector using the affine transformation whose result is forwarded 
and computing the result of the Boolean combination with additional ReLU layers.
For the score predicates $S(i,j)$ we note that they either only depend on $j$ 
or they check whether some binary numbers at positions $i$ and $j$ are equal.
The former can be simulated using an additional preliminary layer that already computes the result of $S(j)$ for every position $j$.
For the latter we use \cref{lem:equality}.
\end{proof}

The proof of the following Lemma can be found in \cref{app:emptiness}.
\begin{lemma}\label{lem:equality}
Given a mask predicate $M$ and tie-breaking function $\tau$,
there is an attention layer using $M$ and $\tau$ that on every sequence $\bm v_1,\dots,\bm v_n \in \{0,1\}^{2d}$
with $\bm v_k = (b_{k,1},1-b_{k,1},\dots,b_{k,d},1-b_{k,d})$ for all $k \in [1,n]$
and for every $i \in [1,n]$ picks attention vector $\bm a_i = \bm v_j$ such that
$b_{i,r} = b_{j,r}$ for all $r \in [1,d]$
if such an unmasked position $j$ exists.
\end{lemma}

We observe that the B-RASP program constructed in the proof of \cref{prop:hardness} 
only uses strict future masking and rightmost tie-breaking.
Thus, the $\EXPSPACE$ lower bound already holds for UHATs that only use strict future masking and rightmost tie-breaking
(similar for strict past masking and leftmost tie-breaking).
\begin{corollary}
The non-emptiness problem for UHATs, where every layer uses strict future masking and rightmost tie-breaking
(resp.\ strict past masking and leftmost tie-breaking), is already $\EXPSPACE$-hard.
\end{corollary}

We now prove the upper bounds in \cref{thm:UHAT-emptiness}.
To this end, we first note that any B-RASP program can be converted in exponential time into an LTL formula
using the construction from \citep{YangCA24}.
In \cref{prop:UHAT-LTL} we prove that the same is true for UHATs,
which improves the doubly exponential construction in \citep{YangCA24} that translates UHATs to B-RASP programs first.
This suffices to prove the exponential-space upper bounds in \cref{thm:UHAT-emptiness} since
non-emptiness of languages given by LTL formulas can be checked in polynomial space \citep{SistlaC85}.

%\begin{proposition}\label{prop:UHAT-BRASP}
%Given a UHAT that recognizes a language $L$, one can construct in polynomial time a B-RASP program recognizing $L$.
%\end{proposition}
%\begin{proof}[Proof sketch]
%We first observe that the values produced by a UHAT are exponentially bounded, 
%since only the result of linear transformations can be forwarded by each layer.
%Here, note that even in an attention layer the dot product is only used to compute the scores, which cannot be forwarded.
%Thus, a polynomial number of bits suffice to encode any value in the UHAT, 
%i.e., the number $B$ in Lemma~22 of \citep{YangCA24} is polynomial.
%We then use the ``smaller version'' of the construction in \citep{YangCA24}.
%Here, we have to show that the score functions can be converted to polynomially sized Boolean expressions
%as opposed to the construction in Lemma~23 of \citep{YangCA24}.
%\end{proof}

For the translation from UHAT to LTL, we first have to make the crucial observation 
that the values occurring during the computation of a UHAT are not ``too large''.

\begin{proposition}\label{lem:UHAT-values}
The values occurring in the computation of a UHAT $\mathcal{T}$ 
can be represented with only a polynomial number of bits in the size of $\mathcal{T}$.
Thus, rational numbers with fixed precision, that is polynomial in the size of $\mathcal{T}$, are sufficient.
\end{proposition}
\begin{proof}
Clearly, ReLU layers do not increase the amount of bits needed since only the maximum with 0 is taken.
Since our UHAT model is defined over rational numbers, 
we can represent each value with a binary number for the numerator and a binary number for the denominator. 
By taking the LCM, we can further assume that the denominators of the numbers in $\emb(\Sigma)$ are all equal.
Note that the LCM of linearly many numbers of linear bit size can be represented with polynomially many bits.
Similarly, for each attention layer we can assume that the denominators of all coefficients of the affine transformation,
whose result is forwarded, are equal.
This ensure that in the input sequence and after each layer the values have the same denominator.
Next we argue that the numerators and denominators of all values that occur in the computation
can be represented with a polynomial number of bits.
The output of an attention layer is an affine transformation involving two input vectors.
Here, the input values are only multiplied with constants, i.e., the values in the output only depend linearly on input values.
Since the denominators of the values after multiplication with coefficients are all equal,
addition does not incur an increase in bit length of the denominators.
Therefore, a repeated application of linearly many attention layers can only lead to numerators and denominators
whose values are at most exponential in the number of layers, i.e., can be represented with polynomially many bits. 
For the score function we observe that the dot product after an affine transformation only involves two input vectors.
Thus, the number of bits needed to represent the result of the score function is still polynomial.
A crucial observation is that the results after applying the score function are not forwarded to the next layer,
which would introduce an exponential increase in size in the number of layers.
\end{proof}

By \cref{lem:UHAT-values}, we can already compute the results of the affine transformations and score functions 
during the construction of the LTL formula.
This means that the LTL formula only has to simulate the position-wise behavior of attention layers, i.e.,
masking and selecting the position of the attention vector,
but not the actual computation of values.
The proof of the following proposition can be found in \cref{app:emptiness}.
\begin{proposition}\label{prop:UHAT-LTL}
Given a UHAT that recognizes a language $L$, one can construct in exponential time an LTL formula recognizing $L$.
\end{proposition}


We remark that if we start with a UHAT, where every attention layer uses strict future masking and leftmost tie-breaking
(resp.\ strict past masking and rightmost tie-breaking),
then the LTL formula constructed in the proof of \cref{prop:UHAT-LTL} only uses the $\past$ (resp.\ $\future$) operator.
It was shown in \citep{SistlaC85} that the non-emptiness problem for the fragments of LTL 
that only allow $\past$ or $\future$ is $\NP$-complete.
Thus, we obtain an improved complexity upper bound for such restricted UHATs.
\begin{corollary}
The non-emptiness problem for UHATs, where every attention layer uses strict future masking and leftmost tie-breaking
(resp.\ strict past masking and rightmost tie-breaking), is in $\NEXP$.
\end{corollary}

Note that it was already shown in \citep{tale} that such restricted UHATs 
are equally expressive as the LTL fragment with only $\past$ (resp.\ $\future$).
However, the construction by \cite{tale} from UHAT to the LTL fragments incurs a doubly exponential blow-up, 
as opposed to our singly exponential translation.
